%%%%%%%%%%%%%%%%%%%%%%%%%%%%%%%%%%%%%%%%%%%%%%%%%%%%%%%%%%%%%%%%%%%%%%
% Lớp 11 | Quan hệ vuông góc trong không gian | Hình lăng trụ đứng. Hình chóp đều. Thể tích
% Dạng 1H8?7-6: Thể tích khối chóp cụt và các khối khác
% Số câu: 4  (sinh tự động bởi tach_cau.py; từ nay thầy sửa trực tiếp file này)
%%%%%%%%%%%%%%%%%%%%%%%%%%%%%%%%%%%%%%%%%%%%%%%%%%%%%%%%%%%%%%%%%%%%%%

%%%%%%%%%%%%%%%% TRẮC NGHIỆM 4 PHƯƠNG ÁN %%%%%%%%%%%%%%%%

% nguồn: Ok_TracNghiemLop11_Lan1
\begin{ex}%[1H8N7-6]%[ma:82c35ee7]
	Cho hình chóp cụt đều $ABC. A'B'C'$.  Tứ giác $AA'B'B$ là hình gì?
	 \choice
		{Hình  vuông}
		{Hình  thoi}
		{\True  Hình thang cân}
		{Hình chữ  nhật}
	\loigiai{ 
	Mỗi mặt bên của hình chóp cụt đều là một  hình thang cân.}
\end{ex}

% nguồn: Ok_TracNghiemLop11_Lan1
\begin{ex}%[1H8N7-6]%[ma:11aefec0]
	Cho khối hộp chữ nhật có 3 kích thước $3;\,4;\,5$.  Thể tích của khối hộp đã cho bằng?
	\choice 
		 {$10$}
		 {$20$}
		{$12$} 
		 {\True $60$}
	\loigiai{ 
		Thể  tích khối hộp chữ nhật $V=a\cdot b \cdot c=3\cdot 4\cdot 5=60$ (đvtt).
	} 
\end{ex}


%%%%%%%%%%%%%%%% TRẢ LỜI NGẮN %%%%%%%%%%%%%%%%

% nguồn: Ok_TraLoiNgan_Lop11_Lan1
\begin{ex}%[1H8V7-6]%[ma:6094b56c]
	Một hình chóp cụt đều ${A B C \cdot A’ B’ C’}$ có  cạnh đáy lớn bằng ${4 a}$, cạnh đáy nhỏ bằng ${2 a}$ và chiều cao của nó bằng ${\dfrac{3 a}{2}}$. Tìm thể tích của khối chóp cụt đều đó với $a=2\sqrt{3}$.
	 \shortans{252}
	\loigiai{ 
		\begin{center}
			\begin{tikzpicture}[scale=1,font=\footnotesize,line join=round,line cap=round,>=stealth]
					 \path
						 (0,0) coordinate (A)
						 ++(-60:3) coordinate (B)
						(6,0) coordinate  (C)
						 ($(B)!1/2!(C)$) coordinate (K)
						($(A)!2/3!(K)$) coordinate (O) 
						($(O)+(0,5)$) coordinate (S) 
						($(S)!.5!(A)$)  coordinate (A')
						 ($(S)!.5!(B)$) coordinate (B')
						($(S)!.5!(C)$)  coordinate (C')
						($(B')!1/2!(C')$) coordinate (J) 
						($(A')!2/3!(J)$) coordinate  (I)
						($(A)!.3!(K)$) coordinate (H); 
					\foreach \i in{A,B,C}{\draw  (\i')--(\i);};
					\draw  (A)--(B)--(C)
						(A')--(B')--(C')--cycle  (A')--(J);
					 \draw[dashed] (I)--(O)
						 (A)--(C)
						(A)--(K) 
						 (A')--(H);
					\pic[draw,angle eccentricity=1.8,angle radius=2mm]{right angle=S--O--K}; 
					\pic[draw,angle  eccentricity=1.8,angle radius=2mm]{right angle=A'--H--O};
					\pic[draw,angle eccentricity=1.8,angle radius=2mm]{right  angle=O--I--J};
					\foreach  \i/\g in {A/180,B/-120,C/0,K/-50,O/-90,A'/90,B'/180,C'/90,J/-50,I/30,H/-90}
					\fill[black] (\i) circle (1pt)+(\g:4mm) node[scale=1]{$\i$}; 
				\end{tikzpicture}
		\end{center}
		Gọi ${O, I}$ theo thứ tự là tâm của đáy lớn ${A B C}$ và đáy bé ${A’ B’ C’;  K, J}$ theo thứ tự là trung điểm của ${B C}$ và ${B’ C’}$.\\
		Ta  có ${h=I O=\dfrac{3 a}{2}}$ là chiều cao của hình chóp cụt đều ${A B C \cdot A’ B’ C’}$.\\
		Diện tích hai đáy hình  chóp cụt đều là\\
		${S_1=S_{\Delta A B C}=\dfrac{(4 a)^2 \sqrt{3}}{4}=4 a^2 \sqrt{3};  S_2=S_{\Delta A’ B’ C’}=\dfrac{(2 a)^2 \sqrt{3}}{4}=a^2 \sqrt{3}
		}$.\\ 
		 Thể tích khối chóp cụt đều là\\
		 $V=\dfrac{1}{3}h\left(S_1+\sqrt{S_1S_2}+S_2\right)$
		$=\dfrac{1}{3}\cdot  \dfrac{3a}{2}\left(4a^2\sqrt{3}+\sqrt{4a^2\sqrt{3}\cdot a^2\sqrt{3}}+a^2\sqrt{3}\right)=\dfrac{7a^3\sqrt{3}}{2}$ (đơn vị thể tích).\\
		Với  $a=2\sqrt{3}$ thì $V=252$.
	 }
\end{ex}


%%%%%%%%%%%%%%%% TỰ LUẬN %%%%%%%%%%%%%%%%

% nguồn: Ok_TuanLuan_Toan11_Lan2
\begin{ex}%[1H8H7-6]%[ma:3f3d5d03]
	 Thể tích khối lăng trụ lục giác đều có tất cả các cạnh bằng $a$.
	\loigiai{ 
	 \immini{
		Ta có $OAB$ là tam giác đều cạnh $a$ nên  $S_{OAB}=\dfrac{\sqrt{3}a^2}{4}$.\\
		 Diện tích lục giác đều $ABCDEF$ là
		\[S_{ABCDEF}=6\cdot S_{OAB}=6\cdot\dfrac{\sqrt{3}a^2}{4}=\dfrac{3\sqrt{3}a^2}{2}.  \]
		Vậy thể khối khối lăng  trụ là
		$V=a\cdot \dfrac{3\sqrt{3}a^2}{2}=\dfrac{3\sqrt{3}a^3}{2}$. 
		}{\begin{tikzpicture}[scale=0.8,font=\footnotesize,line join=round,line cap=round,>=stealth]
		\def\h{3} 
		 \path
			(-1,4)  coordinate (A)
			(.5,5) coordinate  (B)
			 (2,5) coordinate (C)
			(1,4) coordinate (O) 
			($(A)!2!(O)$) coordinate (D) 
			($(B)!2!(O)$) coordinate (E) 
			($(C)!2!(O)$) coordinate  (F)
			(A)+(-90:\h) coordinate  (A')
			 (B)+(-90:\h) coordinate (B')
			(C)+(-90:\h) coordinate  (C')
			(D)+(-90:\h) coordinate (D') 
			(E)+(-90:\h) coordinate  (E')
			(F)+(-90:\h) coordinate (F'); 
		\draw (A)--(B)--(C)--(D)--(E)--(F)--cycle (A)--(A')--(F')--(E')--(D')--(D) 
			(E)--(E') 
			(F)--(F') 
			(A)--(D) 
			(B)--(E) 
			 (C)--(F);
		 \draw[dashed](A')--(B')--(C')--(D')
			 (B)--(B')
			 (C)--(C');
		\foreach \x/\g in  {A/180,B/90,C/90,D/0,E/-40,F/210,O/-90} \fill (\x) circle (0.03)+(\g:3mm) node[scale=.8]{$\x$};
	\end{tikzpicture}
		 }
	 }
\end{ex}
